\centerline{\bf \Huge Матлогика}
\title{Матлогика}
\begin{flushright}
    \scriptsize Существует, так называемая, «Дилемма дикобразов». \\Если они захотят согреть друг друга, то чем тесней они сблизятся,\\ тем сильнее будет боль от уколов игл. То же самое и с людьми...\\
    Рицуко Акаги. Евангелион.
\end{flushright}
\textbf{\Large Взвешивания}

\problem{1}Одна из девяти монет фальшивая, весит легче настоящей. Как определить ФМ за 2 взвешивания?\\

\problem{2}а) Одна из 27 монет фальшивая, она весит легче настоящей. Как определить ФМ за 3 взвешивания?\\
б) Решите ту же задачу для 26 монет. \\

\problem{3}Есть одна золотая, 3 серебряные и 5 бронзовых медалей. Известно, что одна из них фальшивая (весит легче настоящей). Настоящие медали из одного металла весят одинаково, а из различных - нет. Как за 2 взвешивания на чашечных весах без гири найти фальшивую медаль?\\

\problem{4}Есть 5 серебрянных и 4 золотые монеты. Известно, что одна из них фальшивая, а остальные настоящие (настоящая серебряная монета отличается по весу от настоящей золотой). Если ФМ серебряная, то она легче настоящих серебряных монет, а если золотая, то тяжелее настоящих золотых. Как найти ФМ за два взвешивания?\\

\textbf{\Large Высказывания}\\

\problem{1}В день рождения Эльзяты Яна хочет выяснить, сколько той лет. Эвена говорит, что Эльзяте больше 14 лет, а Саша утверждает, что больше 13 лет. Сколько лет Эльзяте, если известно, что ровно одна из них ошиблась? Ответ обоснуйте.    \\

\problem{2}В чашке, стакане, кувшине и банке находятся молоко, лимонад, квас и вода. Известно, что вода и молоко не в чашке; сосуд с лимонадом стоит между кувшином и сосудом с 
\newpage
\noindent
квасом; в банке не лимонад и не вода; стакан стоит около банки и сосуда с молоком. В какой сосуд налита каждая из жидкостей?\\

\problem{3}  Про группу из пяти человек известно, что:\\
Лера на 1 год старше Басанговой,\\
Федор на 2 года старше Лободина,\\
Алтан на 3 года старше Насунова,\\
София на 4 года старше Гиберт,\\
а еще в этой группе есть Юля и Менглинова.\\
Кто старше и на сколько: Юля или Менглинова?\\

\problem{4}Мартышка, Осёл и Козёл затеяли сыграть трио. Уселись чинно в ряд, Мартышка справа. Ударили в смычки, дерут, а толку нет. Поменялись местами, при этом Осёл оказался в центре. А трио всё нейдёт на лад. Пересели ещё раз. При этом оказалось, что каждый из трёх "музыкантов" успел посидеть и слева, и справа, и в центре. Кто где сидел на третий раз?\\

\textbf{\Large Рыцари и Лжецы}\\

\problem{1}Ученики в ЛМШ делятся на тех, кто обучается в 5-6 и 7-8 классах. Известно, что ученики 5-6 классов всегда говорят правду, а 7-8 классы — всегда лгут. Начальник ЭлМШ познакомился с одним из учеников ЛМШ, и они пошли с Пагоды до Дружбы. По дороге они встретили нового ученика ЛМШ. Начальник попросил у своего знакомого ученика узнать, в каком классе обучается новый ученик. Знакомый ученик вернулся и сообщил, что новый ученик сказал, что учится в 5-6. В каком классе учился знакомый ученик — 5-6 или 7-8?\\

\problem{2}Каждый из четырех гномов — Беня, Веня, Женя, Сеня — либо всегда говорит правду, либо всегда врет. Мы услышали такой разговор: Беня — Вене: "ты врун"; Женя — Бене: "сам ты врун"; Сеня — Жене: "да оба они вруны, — (подумав), — впрочем, ты тоже". Кто из них говорит правду?\newpage

\problem{3}Врун всегда лжёт, Хитрец говорит правду или ложь, когда захочет, а Переменчик говорит то правду, то ложь попеременно. Путешественник встретил Вруна, Хитреца и Переменчика, которые знают друг друга. Сможет ли он, задавая им вопросы, выяснить, кто есть кто?\\